% Corte mecánico íntegro: parte_iv.tex:1--54.
% Residencia material del ensamblaje sucesor de 102 capítulos.
% Residencia de realización: capítulo 67.
\part{Differential Geometry and Dimensional Realizations}

\chapter{Curvature and torsion of the published curves}

\section{Planar polar curve}

Let
\[
 \gamma(\theta)=r(\theta)(\cos\theta,\sin\theta).
\]
With primes denoting derivatives with respect to \(\theta\), the velocity and curvature are
\begin{align}
 v(\theta)&=\sqrt{r^2+(r')^2},\\
 \kappa_{\mathrm{F}}(\theta)&=
 \frac{r^2+2(r')^2-rr''}
 {\bigl(r^2+(r')^2\bigr)^{3/2}}.
 \label{espiral:eq:curvatura-polar}
\end{align}
The arc length between \(\theta_0\) and \(\theta_1\) is
\[
 s=\int_{\theta_0}^{\theta_1}
 \sqrt{r^2+(r')^2}\,\dd\theta.
\]
These magnitudes belong to the published curve. They are neither the TRIT regime nor
the radial cocycle.

\section{Specializations}

For the logarithmic spiral \(r=r_0\e^{b\theta}\),
\[
 r'=br,\qquad r''=b^2r,
\]
and
\begin{equation}
 \kappa_{\mathrm{F}}(\theta)
 =\frac{1}{r(\theta)\sqrt{1+b^2}}.
 \label{espiral:eq:curvatura-log}
\end{equation}
The angle \(\psi\) between the tangent vector and the radial direction satisfies
\[
 \tan\psi=\frac{r}{r'}=\frac1b,
\]
and is therefore constant. This is the equiangular property of the logarithmic
spiral.


For the Archimedean spiral \(r=a\theta\),
\[
 \kappa_{\mathrm{F}}(\theta)
 =\frac{\theta^2+2}{a(\theta^2+1)^{3/2}}.
\]
For Fermat \(r=a\sqrt\theta\), the singularity at the origin requires
\(\theta>0\) to be declared. For the hyperbolic and lituus branches, the
angular origin is an asymptote and not a point of the domain.
